\documentclass[10pt,a4paper]{article}
\usepackage[a4paper,margin=1.5cm]{geometry}
\usepackage[utf8]{inputenc}
\usepackage[T1]{fontenc}
\usepackage[brazilian]{babel}
\usepackage{enumitem}
\usepackage[hidelinks]{hyperref}
\usepackage{titlesec}

\setlist[itemize]{leftmargin=1.2em, topsep=2pt, itemsep=2pt, parsep=0pt}
\titleformat{\section}{\large\bfseries\MakeUppercase}{}{0pt}{}[\vspace{-4pt}\hrule]
\titlespacing*{\section}{0pt}{10pt}{6pt}

\pagestyle{empty}

\begin{document}

\begin{center}
{\Large\bfseries Marcelo Vironda Rozanti}\\[2pt]
Engenheiro de Software Backend Sênior | Java | Kotlin | Spring Boot | Pix | Microsserviços\\[2pt]
São Paulo, Brasil | mvrozanti@hotmail.com | linkedin.com/in/mvrozanti | github.com/mvrozanti | mvrozanti.github.io
\end{center}

\section{Resumo profissional}
Engenheiro de Software Backend Sênior com mais de 10 anos de experiência no desenho, construção e operação de sistemas backend de produção para bancos, pagamentos e serviços financeiros regulados. Especialista em Java e Kotlin sobre o ecossistema Spring / Spring Boot. Hands-on em PostgreSQL, MongoDB, MySQL, Oracle, MSSQL, Redis. Experiência com infraestrutura AWS, containerização Docker, CI/CD (GitLab, Harness, Travis), Terraform (IaC), observabilidade (Splunk, Grafana, Dynatrace), segurança de aplicação (Fortify, Sonar, CyberArk). Expertise em fintech brasileira: pagamento instantâneo Pix, automação fiscal SPED / eSocial / EFD-Reinf, conformidade com a LGPD, decisão de crédito e registro de recebíveis (Duplicatas Escriturais via B3 / CERC). Entrega cada vez mais assistida por IA: Retrieval-Augmented Generation (RAG), avaliação de LLM, agentic coding, método BMAD e MCP. Português nativo; inglês fluente; espanhol intermediário.

\section{Habilidades técnicas}
\begin{itemize}
\item \textbf{Linguagens}: Java (8, 11, 15, 16, 17, 21), Kotlin, Python, Shell / Bash, C, C++
\item \textbf{Frameworks}: Spring, Spring Boot, Spring Security, Spring Data, MyBatis, Javalin, Koin, JSF, PrimeFaces, Java EE
\item \textbf{Bancos de dados}: PostgreSQL, Aurora PostgreSQL, MongoDB, MySQL, Oracle SQL Server, Microsoft SQL Server, Redis
\item \textbf{Cloud \& DevOps}: AWS (SQS, SNS, S3, EC2, Lambda), Docker, Terraform / HCL, Harness, GitLab CI/CD, Travis CI, Maven, Ant
\item \textbf{Mensageria \& auth}: RabbitMQ, Keycloak
\item \textbf{IA / LLM}: Retrieval-Augmented Generation (RAG), avaliação de LLM, agentic coding, método BMAD, MCP, engenharia de prompt, saída estruturada / guardrails
\item \textbf{Observabilidade \& segurança}: Splunk, Grafana, Dynatrace, Graylog, LogDNA, Fortify, SonarQube, CyberArk
\item \textbf{Sistema operacional}: Linux (uso diário desde 2009)
\item \textbf{Metodologias}: Domain-Driven Design (DDD), Test-Driven Development (TDD), DevSecOps, CI/CD, Agile, DTAP
\item \textbf{Domínios}: Bancário, pagamentos, Pix, POS / presencial, decisão de crédito, automação fiscal, LGPD, registro de recebíveis (Duplicatas Escriturais, B3 / CERC)
\end{itemize}

\section{Experiência profissional}

\noindent\textbf{SRM Asset} -- Engenheiro de Software Sênior \hfill Março 2026 -- Atual\\
São Paulo, Brasil
\begin{itemize}
\item Integrações backend com as registradoras B3 e CERC para Duplicatas Escriturais (recebíveis escriturais), trilho regulado de registro de recebíveis no Brasil
\item Integração com as APIs da QI Tech (banking- e credit-as-a-service)
\item Entrega acelerada por IA: Retrieval-Augmented Generation (RAG) sobre documentação regulatória e de APIs, método ágil-IA BMAD, boas práticas de LLM (avaliações, saída estruturada, guardrails, MCP, engenharia de prompt)
\item Java (8-21), Spring, Aurora PostgreSQL, AWS, Docker, MyBatis, RabbitMQ, Keycloak, Terraform / HCL
\end{itemize}

\noindent\textbf{C6 Bank} -- Engenheiro de Software Sênior \hfill Abril 2025 -- Março 2026\\
São Paulo, Brasil
\begin{itemize}
\item Desenha, implementa e mantém sistemas backend de alta disponibilidade para concessão de crédito em um dos maiores bancos digitais do Brasil
\item Domain-Driven Design nos domínios de decisão de crédito em Kotlin e Spring
\item Foco em segurança da informação: gestão de segredos, princípio do menor privilégio, CI/CD endurecido via Harness
\item Observabilidade com Splunk e Grafana; camada de dados em MongoDB e PostgreSQL; deploys containerizados via Docker
\end{itemize}

\noindent\textbf{Fiserv} -- Engenheiro de Software \hfill Março 2024 -- Abril 2025\\
São Paulo, Brasil
\begin{itemize}
\item Entrega de serviços backend para a plataforma POS de pagamentos presenciais
\item Responsável pela postura de prevenção e recuperação de desastres em produção
\item Liderou adoção de DevSecOps: gates Fortify e SonarQube em CI, CyberArk para segredos, pipelines de release endurecidos
\item Operação de produção com Splunk e Dynatrace; camada de dados Oracle SQL Server e Redis
\item Documentação técnica e processual que suporta rotação de time
\end{itemize}

\noindent\textbf{Upwork} -- Desenvolvedor Backend (Freelance) \hfill Setembro 2022 -- Fevereiro 2024\\
Remoto
\begin{itemize}
\item Contratado autônomo para clientes internacionais: serviços backend e pipelines de extração de dados / scraping e automação web sobre fontes diversas
\item Comunicação com clientes em inglês: levantamento, escopo e entrega
\item Python (lxml, BeautifulSoup, pup), navegadores headless (Chromium, PhantomJS), Bash
\end{itemize}

\noindent\textbf{Conta Azul} -- Desenvolvedor Backend \hfill Setembro 2021 -- Agosto 2022\\
São Paulo, Brasil
\begin{itemize}
\item Construção de motor de automação de obrigações fiscais brasileiras (SPED, eSocial, EFD-Reinf)
\item Processamento periódico e orientado a eventos sobre primitivos nativos AWS (SQS, S3)
\item Infraestrutura como código via Terraform; camada de dados PostgreSQL
\item Java 8/11/15/16 com Spring; observabilidade via LogDNA
\end{itemize}

\noindent\textbf{Grão} -- Desenvolvedor Backend \hfill Dezembro 2020 -- Setembro 2021\\
São Paulo, Brasil
\begin{itemize}
\item Implementação da integração com o sistema brasileiro de pagamento instantâneo (Pix) em Kotlin, do kickoff ao lançamento em produção
\item Modelagem de esquema e entidades focada em performance
\item Design de APIs REST escaláveis; automação de QA no loop de entrega
\item Ambientes de teste e produção containerizados; Java EE, Kotlin, MySQL, Spring, Feign Client, Firebase, GitLab CI/CD, Graylog, Newman/Postman, Tomcat, Svelte
\end{itemize}

\noindent\textbf{Raia Drogasil} -- Desenvolvedor de Software \hfill Janeiro 2020 -- Dezembro 2020\\
São Paulo, Brasil
\begin{itemize}
\item Adequação do software de Operações de Venda à Lei Geral de Proteção de Dados (LGPD)
\item Integração de leitores biométricos e impressoras térmicas no stack de loja
\item Modelagem diária em Oracle SQL; Java EE, Java Applets, JSF/PrimeFaces, JBoss, WebLogic, Spring
\item Primeiro contato com DTAP (Development-Test-Acceptance-Production); TDD
\end{itemize}

\noindent\textbf{Sys4Bank} -- Desenvolvedor de Software \hfill Março 2016 -- Janeiro 2020\\
São Paulo, Brasil
\begin{itemize}
\item Desenvolvimento backend para infraestrutura bancária: fluxos Corporativo, Financeiro, Crédito e Conta
\item Introdução de containerização e pipelines Travis CI em time vindo de deploy manual
\item Modelagem diária em Microsoft SQL Server
\item Java, Spring Boot, MSSQL, Maven, Ant, Bash, Linux
\end{itemize}

\section{Formação}
\noindent\textbf{Bacharelado em Ciência da Computação}, Universidade Presbiteriana Mackenzie, São Paulo, Brasil -- 2020\\
Monografia: \textit{Três experimentos com escalonamento assíncrono de atualização por prioridade de vizinhança em autômatos celulares elementares}

\section{Open source}
\begin{itemize}
\item \textbf{RAT-via-Telegram} -- 600+ estrelas -- Windows RAT via Telegram (Python) -- github.com/mvrozanti/RAT-via-Telegram
\item \textbf{dte} -- 300+ estrelas -- Linguagem de processamento de data e hora -- github.com/mvrozanti/dte
\item \textbf{yawc} -- Yet Another WhatsApp Client -- github.com/mvrozanti/yawc
\item \textbf{dotty} -- Versionamento e sincronização de dotfiles -- github.com/mvrozanti/dotty
\end{itemize}

\section{Idiomas}
Português (nativo), Inglês (fluente), Espanhol (intermediário)

\end{document}
